\section{Exhaustive finite certification}\label{finite:sec:catalog}
The argument in this section is computer assisted.  The unrestricted
bounds proved earlier reduce the categorical problem to finitely many
integral tensors. We enumerate every ordinary fusion ring with $D<1000$
and retain an explicit certificate for each input. Independently, a complete
signed cover through multiplicity $31$ supplies all possible sphericalization
quotients needed by the general-case exclusions. Both programs and both data
sets are included; neither uses an external census or a pivotal classification
as an oracle.

\subsection{Tensor conventions and exact elimination}\label{finite:sec:census}
Every rank-four fusion ring is commutative by
Lemma~\ref{core:lem:rank-four-algebra}.
For a self-dual basis $1,x,y,z$, reciprocity gives
\[
N_x=\begin{pmatrix}0&1&0&0\\1&a&b&c\\0&b&d&e\\0&c&e&f\end{pmatrix},\quad
N_y=\begin{pmatrix}0&0&1&0\\0&b&d&e\\1&d&g&h\\0&e&h&i\end{pmatrix},\quad
N_z=\begin{pmatrix}0&0&0&1\\0&c&e&f\\0&e&h&i\\1&f&i&j\end{pmatrix}.
\]
These formulas also define signed tables, with integer coefficients.
Associativity is equivalent to the following six equations:
\begin{align}
 ad-b^2+bg+ch-d^2-e^2+1&=0,\label{finite:eq:census1}\\
 ae-bc+bh+ci-de-ef&=0,\label{finite:eq:census2}\\
 be-cd+dh-eg+ei-fh&=0,\label{finite:eq:census3}\\
 af+bi-c^2+cj-e^2-f^2+1&=0,\label{finite:eq:census4}\\
 bf-ce+di-eh+ej-fi&=0,\label{finite:eq:census5}\\
 df-e^2+gi-h^2+hj-i^2+1&=0.\label{finite:eq:census6}
\end{align}
This follows either by comparing the three pairwise matrix commutators
or by direct substitution in the associativity identities.  If $b,c>0$, put
\begin{gather*}
 H_0=e(d+f)+bc,\quad G_0=b^2+d^2+e^2-1,\quad J_0=c^2+e^2+f^2-1,\\
 A=e(b^2-c^2)+bc(f-d),\qquad C=e(bf-ce),\\
 B_0=b^3e-b^2cd+b(d-f)H_0-ebG_0+ecH_0.
\end{gather*}
Then
\begin{equation}\label{finite:eq:census-linear}
 Ca+Ai+B_0=0,
\end{equation}
with
\[
 h=(H_0-ae-ci)/b,\qquad g=(G_0-ad-ch)/b,
 \qquad j=(J_0-af-bi)/c.
\]
Divisibility, bounds and all six original equations are checked after
this elimination.  In particular the linear equation is used only as a
necessary condition and does not replace the original equations.
\subsection{The complete ordinary census}\label{finite:sec:ordinary-census}
The ordinary census uses $M=31$ and the strict dimension condition $D<1000$.
For the displayed matrices, put $B=\sum_i N_iN_i^{\mathsf T}$.
The positive dimension vector is an eigenvector of this symmetric
nonnegative matrix with eigenvalue $D$. Thus $D$ is its largest eigenvalue,
and
\[
 D<1000\quad\Longleftrightarrow\quad 1000I-B\text{ is positive definite}.
\]
The generator and checker test its leading principal minors by exact integer
arithmetic. The boundary $D=1000$ is excluded, not rounded.
Appendix~\ref{app:ordinary-census} gives every elimination branch and proves
that the resulting census contains each ring below the bound exactly once.
It has $\CensusTotal$ classes, split into $\CensusSelfDual$ self-dual classes
and $\CensusNonSelfDual$ classes with a dual pair.

The local source and data are in \texttt{unrestricted/supplementary/census}.
The full tensors, their dualities, the ten- or six-parameter presentations,
canonical fingerprints, exact dimension certificates, and counts at each
multiplicity are supplied. The independently implemented Python tensor
checks verify all fusion-ring identities and all unit-preserving basis
permutations. Completeness rests on the exhaustive branch proof, rather
than only on agreement of totals.

\subsection{A shared exhaustive signed cover}\label{finite:sec:shared-signed-cover}
The signed reduction permits an arithmetic exclusion to be shared among
all unsigned rings compatible with a signed table. The following
construction makes that shared search exhaustive.

\begin{proposition}[A shared exhaustive signed cover]
\label{finite:prop:global-signed-cover}
Fix $M\geq1$. Consider the integral self-dual rank-four signed tables
with the prescribed unit and reciprocity, parameter tuple
$(a,b,c,d,e,f,g,h,i,j)$, and all ten coefficients of absolute value
at most $M$. The six associativity equations define a finite set
preserved by the group of nonunit basis permutations and basis-sign
changes, of order $48$.

Every orbit has a representative satisfying
\[
 b,c,d\geq0,\qquad
 b=\min\{|b|,|c|,|d|,|f|,|h|,|i|\},
 \qquad e\geq0\ \text{if }b=c=0.
\]
Enumerating all such representatives and retaining only the
lexicographically least preferred tuple in each of their $48$
transforms gives exactly one table per orbit.

For every self-dual unsigned fusion ring $R$ of multiplicity at
most $M$, all its possible sphericalization quotients are covered
by the transforms of these representatives that satisfy its signed
coefficient bounds, parities, and adjoint-parity conditions.
A representative failing the formal-codegree $d$-number criterion
or character cyclotomicity can be discarded for every $R$
simultaneously.
\end{proposition}
\begin{proof}
The permutation group is transitive on the six mixed positions
$b,c,d,f,h,i$. Place a position of minimum absolute value at $b$.
Under a basis-sign change $(s_1,s_2,s_3)$, the signs multiplying
$b,c,d$ are respectively $s_2,s_3,s_1$. Hence these three entries
can be made nonnegative independently, without changing their
absolute values. If $b=c=0$, the remaining choices of $s_2,s_3$
also make $e$ nonnegative while preserving $d\geq0$.
Thus the preferred set meets every orbit. Its least element in
each finite orbit is unique; checking all $48$ transforms selects
precisely that element.

For completeness, the preferred set is exhausted by the following
disjoint branches of the associativity equations. If $b>0$, then
$c,d\geq b$ and $|f|,|h|,|i|\geq b$; the two-coefficient
Diophantine relation \eqref{finite:eq:census-linear}, with the polynomial
coefficients from the preceding subsection, holds for signed
parameters as well. Exhaust the bounded choices of $b,c,d,e,f$
and solve that relation for all bounded integer pairs $a,i$ by
exact greatest-common-divisor progressions. If one coefficient
vanishes, exhaust its free variable; if both vanish, retain the
entire bounded product precisely when the constant term vanishes.
The formulas in that subsection then determine $h,g,j$. Impose the
signed bounds and the preferred-representative conditions on all
coordinates, and check the original six equations.
If $b=0<c$ and $e\ne0$, the equations determine
\begin{align*}
 h&=\frac{d^2+e^2-ad-1}{c},&
 i&=\frac{e(d+f-a)}c,\\
 j&=\frac{c^2+e^2+f^2-af-1}{c},&
 g&=\frac{h(d-f)+ei-cd}{e}.
\end{align*}
The divisibility for $i$ gives
$a\equiv d+f\pmod{c/\gcd(c,e)}$; all bounded choices in this
progression are checked. If $b=0<c$, $e=0$ and $d>0$, then
\[
 \gcd(c,d)=1,\quad i=0,\quad
 h=\frac{d^2-ad-1}{c}\ne0,\quad f=d-\frac{cd}{h},
\]
and $j$ is given by the preceding formula, while $g$ is enumerated.
The remaining case $b=0<c$, $e=d=0$ forces
$(c,h,f,i,j)=(1,-1,0,0,0)$ and leaves $a,g$ free.

If $b=c=0<e$, then $a=d+f$ and $df=e^2-1$;
the equations determine $g,j$ from $d,e,f,h,i$, with the case
$d=0$ retained. Finally, $b=c=e=0$ forces $d=1,a=0$ and
$f\in\{1,-1\}$. For $f=1$ the remaining equation is
\[
 gi+hj=h^2+i^2-2,
\]
including its $h=0$ branch. For $f=-1$ one has $h=i=0$
and $g,j$ are free. Coordinates whose signs are not fixed by the
preferred convention retain both signs and the original size
bounds, and every emitted tuple is checked
against all six equations. This proves exhaustive enumeration
of the preferred set, including the zero denominators.

Any sphericalization quotient for $R$ lies in the bounded signed
set by Lemma~\ref{core:lem:family-induction}. One of the $48$ transforms
of its global representative is the quotient in the basis of $R$,
so testing all compatible transforms loses no categorical
possibility. Formal codegrees are unchanged by permutations and
basis signs, and character values change only by permutation and
sign. Their $d$-number and cyclotomicity conditions are therefore
orbit-invariant. Both conditions are necessary for a signed
sphericalization quotient by Theorem~\ref{finite:cor:joint-signed-conditions},
which proves the common exclusions.
\end{proof}

\begin{remark}[Certificate reuse]
The global signed census is regenerated exhaustively during the
verification workflow; its completeness is not inferred from a
cached list. Universally excluded orbits receive a checked common
arithmetic rejection. Every other orbit is expanded into all
distinct transforms compatible with each unsigned ring, and every
such candidate receives the corresponding signed-branch check.
Sharing these computations changes their allocation among records,
not the necessary criteria or the scope of the classification.
\end{remark}

\subsection{The positive-orbit bridge}
\begin{proposition}\label{finite:prop:positive-bridge}
Let $\mathcal T_M$ be the complete signed cover after the common
$d$-number and cyclotomicity exclusions.  Every self-dual rank-four
Grothendieck ring of multiplicity at most $M$ is isomorphic to a
nonnegative gauge transform of a table in $\mathcal T_M$.
Consequently, certificates for all such nonnegative orbits suffice for the
self-dual classification. The ordinary census is a separate exact input list
for the successive-filter ledger; its completeness is not assumed in this
signed-cover argument.
\end{proposition}
\begin{proof}
The ordinary fusion table itself is an integral associative signed table
with the stated coefficient bound.  Its orbit occurs in the exhaustive
cover, and its formal codegrees and character values satisfy the common
necessary arithmetic conditions.  Thus its representative belongs to
$\mathcal T_M$.  Applying all eight sign gauges and all six nonunit
permutations recovers its ordinary table.

A sign change taking a nonnegative table to another nonnegative table
fixes every nonzero coefficient: the sign multiplying such a coefficient
must be $+1$.  Thus a signed orbit contains at most one nonnegative orbit
modulo permutations.  Equality between the list of these positive orbits
and the list of certificate inputs is an exact finite coverage check.
\end{proof}

Theorem~\ref{thm:uniform-global} reduces the classification to $M=31$.
The fresh exhaustive signed enumeration has $\VerifyAllSigned$ orbits.
The formal-codegree $d$-number condition rejects $\VerifyDnumberRejected$,
and cyclotomicity rejects a further $\VerifyCyclotomicRejected$, leaving
$\VerifyRetainedSigned$. Exactly $\VerifyPositiveRecords$ retained orbits
have a nonnegative transform; each has one checked certificate. This
signed-cover scope is all multiplicities at most $31$, without the ordinary
census's additional condition $D<1000$. Its counts are therefore distinct
from the fusion-ring counts in Table~\ref{finite:tab:filters}.

The supplementary certificate bank keeps the broader signed-cover proof
intact. The full-census ledger tests rank-two subrings and universal
arithmetic directly. Every later self-dual witness carries an explicit basis
isomorphism to its certificate input. A table index, digest, or stored status
is never an exclusion criterion.

The automatic-sphericality branch checks its structural hypothesis
before using either Theorem~\ref{inv:thm:rankfour-invertible-spherical}
or Corollary~\ref{core:cor:multiplicity-one-spherical}. It then checks
the supplied spherical obstruction, including any required
Frobenius--Perron Galois-orbit condition.
The arithmetic and divisibility inputs are stated in
Proposition~\ref{finite:prop:standard}.  The induction and moment tests
are Propositions~\ref{finite:prop:low-certificate},
\ref{finite:prop:optimized-moment} and
\ref{finite:prop:spherical-codegree-moment}; the quadratic root-length
branch is Proposition~\ref{finite:prop:quadratic-family-n3-n4}, including
its conductor-$40$ boundary.  Signed tests use
Lemma~\ref{core:lem:family-induction} and
Theorem~\ref{finite:cor:joint-signed-conditions}, together with the norm and
field obstructions in Lemmas~\ref{finite:lem:norm}
and~\ref{finite:lem:field-prime}.
A reduction to pseudo-unitarity is made only after checking weak
integrality, the two-dimension hypothesis, or a spherical dimension
character in the Frobenius--Perron Galois orbit.  No pivotality assumption
is imposed on the original input category.

For clarity, two different enumerations are verified.  Every compatible
signed table is covered, either by the global cover or by the indicated
complete bounded local search.  For every low-induction certificate,
every bounded integral restriction column and every admissible combination
of columns is regenerated.  In particular, neither a stored classification
flag nor a claimed exhaustive list is accepted without checking its witness
and its coverage.  Character fields are represented by squarefree
polynomials over $\Q$; supplied root lists are checked by multiplying
linear factors.  Algebraic signs use rational isolating intervals.

\begin{theorem}[Computer-assisted self-dual classification]
\label{finite:thm:selfdual-classification}
The self-dual rank-four complex Grothendieck rings of multiplicity at most
$31$ are exactly the nine self-dual rings in Table~\ref{finite:tab:models}.
\end{theorem}
\begin{proof}
Propositions~\ref{finite:prop:global-signed-cover}
and~\ref{finite:prop:positive-bridge} reduce the question to the
$\VerifyPositiveRecords$ positive orbits above.  The exact verifier checks a based-ring
isomorphism for each of the nine model records and accepts a necessary-
condition contradiction for each of the other $\VerifyExcludedRecords$ records.
The mathematical implications are precisely those listed above.
All nine surviving tensors are realized by the cited unitary categories.
\end{proof}

\subsection{The non-self-dual case}
\begin{theorem}\label{finite:thm:nonselfdual-classification}
There are exactly six non-self-dual rank-four complex Grothendieck rings:
the six non-self-dual entries in Table~\ref{finite:tab:models}.
This assertion does not require finite enumeration.
\end{theorem}
\begin{proof}
Corollary~\ref{ns:cor:pseudo-unitary-conjugate} permits the
pseudo-unitary reduction without any initial pivotality hypothesis.
Larson's pseudo-unitary candidate theorem \cite[Theorem~1.1]{Larson}
therefore applies.  It gives the six stated rings and the one additional
candidate $S(2)$.  Theorem~\ref{finite:thm:nonselfdual-quadratic} excludes
$S(2)$ (indeed every $S(m)$ with $m\geq2$).  The six remaining rings have
the unitary realizations in Table~\ref{finite:tab:models}.
\end{proof}

\subsection{Successive filters and complete coverage}\label{finite:sec:successive-filters}
The full-census ledger checks all $\CensusTotal$ input rings and records
$\CensusExcluded$ exclusions and $\CensusModels$ models. Its stages are
shown in Table~\ref{finite:tab:filters}. Each ring is counted exactly once,
at its recorded rejection stage. These are counts for the implemented
ordered certificate pipeline; they do not assert that no earlier criterion
could also exclude a ring assigned to a later stage.

% Generated by supplementary/check_manuscript.py from the checked ledger.
\begin{table}[htbp]
\centering\small
\begin{tabular}{@{}lrr@{}}
\toprule
Stage & Excluded & Remaining \\
\midrule
Complete census, $D<1000$ & --- & 1\,692 \\
Rank-two subring obstruction & 375 & 1\,317 \\
Formal-codegree $d$-numbers & 1\,084 & 233 \\
Character cyclotomicity & 2 & 231 \\
Non-self-dual reduction & 21 & 210 \\
Two-dimension and weak-integral reductions & 109 & 101 \\
Subring dimension divisibility & 3 & 98 \\
Involution moment obstructions & 10 & 88 \\
Automatic sphericality and spherical tests & 6 & 82 \\
Signed characters and central induction & 67 & 15 \\
\bottomrule
\end{tabular}
\caption{Successive exact filters for the complete ordinary fusion-ring
census. Each ring is counted once, at its recorded rejection stage.
The final stage covers every compatible signed quotient and every
character or induction-column choice required by its certificate.}
\label{finite:tab:filters}
\end{table}


The rank-two and universal arithmetic stages apply
Propositions~\ref{finite:prop:ranktwo} and~\ref{finite:prop:standard}.
For a non-self-dual input, the verifier checks a based isomorphism with
one of Larson's seven candidates or its absence from that list; the
extra candidate $S(2)$ is excluded by
Theorem~\ref{finite:thm:nonselfdual-quadratic}.
The reduction to the pseudo-unitary candidate theorem is justified by
Corollary~\ref{ns:cor:pseudo-unitary-conjugate}, not imposed on the input.
The remaining stages replay precisely the necessary conditions described
above, with all reduction hypotheses and basis changes checked.
The final signed stage covers all compatible quotients and all required
character and restriction-column choices, rather than a preferred lift.

\subsection{Models and exact dimensions}
Table~\ref{finite:tab:models} supplies one unitary realization and a
construction reference for every surviving based ring. Here
$\operatorname{Ad}(\mathrm{SU}(2)_k)$ denotes the adjoint subcategory,
$\mathcal H_1$ is the rank-four Haagerup even part, and a near-group ring
of type $(C_3,m)$ satisfies
$\rho^2=\sum_{g\in C_3}g+m\rho$.
The six non-self-dual models are
$\mathrm{Vec}_{C_4}$, the Tambara--Yamagami category on $C_3$,
the principal even part of $S'$, $\operatorname{Rep}(A_4)$, and the
two near-group categories of parameters three and six.

The supplementary model data contain the fifteen complete multiplication
tensors and positive dimension vectors in their actual bases. Exact
multiplication identities identify each vector as the Frobenius--Perron
character. Independent reconstructions from the group, product, adjoint,
near-group, and exceptional fusion rules check all fifteen based-ring
identifications. The polynomials for $\delta_7$ and $\delta_E$ in the table
are checked as factors of the corresponding Casimir characteristic
polynomials; positivity identifies their largest roots. The construction
references prove existence. Sphericality of \emph{every} categorification
is proved separately below.

\begin{corollary}[Computer-assisted classification and sphericality]
\label{finite:cor:classification}
The rings in Table~\ref{finite:tab:models} are exactly the Grothendieck
rings of rank-four complex fusion categories.  Every such category
admits a spherical structure.  Each listed ring has a unitary realization.
\end{corollary}
\begin{proof}
The unrestricted bound of Theorem~\ref{thm:uniform-global} places every
categorifiable ring within the verified multiplicity range.  Apply
Theorems~\ref{finite:thm:selfdual-classification}
and~\ref{finite:thm:nonselfdual-classification}.

For sphericality, the six self-dual multiplicity-one entries are
covered by Corollary~\ref{core:cor:multiplicity-one-spherical}.
The $2D2$ ring has a nontrivial invertible basis element, so
Theorem~\ref{inv:thm:rankfour-invertible-spherical} applies.
The $E_8$ ring is the Fibonacci normal form $(a,b)=(1,1)$, covered
by Proposition~\ref{fib:prop:small-spherical}. The Haagerup ring
is covered by Theorem~\ref{haag:thm:automatic-spherical}, and the
six non-self-dual rings by Theorem~\ref{ns:thm:ns-spherical}.
These arguments apply to every categorification, independently of the
chosen unitary realizations.
\end{proof}

\begin{corollary}[Sharp bounds]\label{finite:cor:sharp-bounds}
Every rank-four complex fusion category satisfies
\[
\operatorname{FPdim}(\C)\leq24+12\sqrt3,
\qquad \max_{i,j,k}N_{ij}^k\leq6.
\]
Both bounds are attained, and equality in either bound holds precisely
for the near-group ring $\operatorname{NG}(C_3,6)$.
\end{corollary}
\begin{proof}
Read the multiplicities and dimensions from Table~\ref{finite:tab:models}.
Every other row has $\mu\le3$ and $D<44$. The radical comparisons are
immediate. For $\delta_7$ and $\delta_E$, substituting $x=44+u$ in their
defining polynomials gives strictly positive coefficients, so neither has
a root at least $44$. Since $44<24+12\sqrt3$, the maximal entry is unique.
The near-group unitary realization proves sharpness of both bounds.
\end{proof}

\subsection{Reproducibility and the scope of the computation}
All code and data are supplied locally. Within \texttt{supplements/rank4}, the directory \texttt{unrestricted/supplementary/census} contains the complete ordinary
fusion-ring enumeration and data. The general-case classification source,
signed cover, witness bank, full-census ledger, and replay reports are in
\nolinkurl{unrestricted/supplementary/classification}. No external census, pivotal package,
or external classification status is read. From the repository root, run
\begin{verbatim}
python3 supplements/rank4/unrestricted/supplementary/check.py
\end{verbatim}
This regenerates the ordinary census, checks every full tensor and exact
dimension bound, regenerates the full-census ledger, independently
regenerates the signed cover, and replays all certificates. It then verifies
all model dimensions and reconstructs both manuscript tables from the
checked data. The requirements are Python, SymPy, and a C++17 compiler
with OpenMP. The optional \texttt{python-flint} backend accelerates the
same exact arithmetic; it is not required.

The \texttt{--deep} flag also removes both dimension-pruning steps from
the ordinary enumerator, applies a separately implemented exact dimension
test, and checks an exhaustive small parameter box. The unpruned search
has $16065$ classes at multiplicity at most $31$; its exact $D<1000$
filter recovers precisely the same $\CensusTotal$ inputs. At multiplicity
two the full box of $3^{10}+3^6$ assignments gives the same $27$ classes.
These comparisons supplement the proof of coverage, with their shared
components stated in the supplementary guide.

The delivered fresh replay checks $\CensusExcluded$ exclusions and all
fifteen model identifications, with no unresolved records. Its signed
subroutine also checks the broader $\VerifyPositiveRecords$-record bank.
The optional local Python signed-search engine is provided as a separate
replay mode. Cached-cover and ledger-only modes are explicitly conditional
and are not used by the complete command above. Character signs are decided
by rational isolating intervals; numerical approximations never determine
an exclusion. These are exact finite checks, not a proof-assistant
formalization of the categorical arguments.
The finite computation verifies the explicitly bounded arithmetic and
coverage assertions. The categorical arguments and the cited category
constructions remain mathematical inputs; the research notice on the
first page applies to both parts of the proof.

